\section{Reproducibility and Engineering Cost}
\label{sec:repro}

\subsection{Software and infrastructure}

Both relational and vector retrieval run on PostgreSQL 17 with pgvector 0.8.2
\citep{pgvector2024}, using an HNSW index over \texttt{vector\_cosine\_ops}
\citep{malkov2018hnsw} alongside \texttt{pg\_trgm} and a GIN \texttt{tsvector}
index. \texttt{hnsw.iterative\_scan} is set to \texttt{relaxed\_order} for reasons
given below. The graph is Neo4j 5.26 Community with APOC \citep{neo4j2024}.
Embeddings come from \texttt{all-MiniLM-L6-v2} at 384 dimensions with a 256
word-piece limit, chunked at 1{,}400 characters with 200 of overlap, batch size
256. Generated SQL and Cypher execute under a 45-second statement timeout.

The deployment is a single application host --- an AWS \texttt{m7g.xlarge}, 4 vCPU
and 16\,GB, running the graph, the API and the gateway --- with PostgreSQL on a
\texttt{db.t4g.medium} instance and CloudFront terminating TLS. There is no NAT
gateway; object-store access goes through a VPC endpoint. \Cref{tab:cost} gives
the monthly cost at roughly \$165. The embedding pass was run on a temporarily
resized 16-vCPU host at 97 chunks/s and the instance was then scaled back.

Ingestion takes about 252 minutes end to end (\Cref{fig:pipeline}), of which
document fetch and embedding are about 82\%. The fetch component is bounded
externally by EDGAR's ten-requests-per-second fair-access policy
\citep{sec2024edgar} and cannot be parallelized away.

\subsection{Auditing the harness}
\label{sec:harness}

Building the measurement code surfaced defects that, unnoticed, would each have
produced a publishable-looking but incorrect result: a pgvector post-filter that
returned 0 of 8 rows while 164 matched, so the vector arm was briefly scored on
an empty retriever; an answer-token cap that truncated answers mid-citation,
silently converting passes into uncited; a global token counter that attributed
every thread's cost to whichever arm read it first. These were found and fixed
before the reported numbers were generated, so they are development history
rather than caveats on the results, and we do not enumerate them here.

Two are worth reporting, because they did reach published numbers.

\textbf{A ground-truth error, found after the first analysis.} One
cross-document question accepted a single expected parent where the loaded data
contains three filers that genuinely list the subsidiary. All three arms were
naming a real parent and being scored wrong. Because every model answer is
stored, correcting the ground truth and re-scoring cost no additional inference;
it moved text-to-SQL from 28 to 29 passing attempts, vector RAG from 6 to 9 and
GraphRAG from 44 to 45. That we found it after believing the results were final
is direct evidence that others may remain.

\textbf{An asymmetry that penalised the graph, and what fixing it changed.}
The more instructive defect is one we initially reported as a finding. Both
stores keep a normalised name column, lowercased with corporate suffixes
stripped, so \emph{AEP Texas Inc.}\ is stored as \texttt{aep texas}. Neither
schema description said so. Text-to-SQL was unaffected because it matched the
raw column with \texttt{ILIKE}; the graph arm wrote an exact property match
against \texttt{nameNormalized} using the name as it appeared in the question,
retrieved zero rows, and refused. We first published that as a category loss for
structural retrieval.

It is not one. It is the same class of defect as an asymmetry we had already
found and fixed in the opposite direction --- an earlier iteration documented a
name convention to the graph but not to SQL, making text-to-SQL fail for a
reason unrelated to its architecture. Having recorded that, we should have
checked for its mirror image and did not.

We therefore documented the normalisation identically in both schema
descriptions, changing no retrieval code and no question, and re-ran all 180
attempts. \Cref{tab:beforeafter} reports both measurements. We present the
repaired numbers as the paper's result and the original alongside them, because
a fix that moves a headline in the author's favour should be visible rather than
quietly absorbed.

The general lesson is not that our harness had bugs. It is that a benchmark
harness is itself an instrument requiring calibration, and that a list of
defects already found is only useful if it is read as a checklist for the ones
not yet found.

\subsection{Artifacts}

The ingestion pipeline, the three retrieval arms, the scoring harness, the
question set with its ground truth, and the stored answers for all 180 attempts
are retained, which is what made the re-scoring in item 8 possible without new
inference. Every table and figure in this paper is generated programmatically from
the stored results rather than transcribed, so the numbers in the prose and the
numbers in the tables cannot drift apart.
